\chapter{Open Questions and a Research Programme}
\label{ch:offen}

% Source: conjectures.md, as of version 0.7.

Before the final chapter locates the theory on the map of
mathematics, an honest stop along the way: not what has been
proved -- the preceding chapters have done that -- but what is
open, and in what order it is worth attempting.

\section{The most immediate questions}

\paragraph{The digit patterns of the factorials.} The digit
formula for the shell jumpers has shifted the question: no longer
``how many jumpers are there?'', but ``how does $n!$ write itself
in its own shell?''. First patterns are in: the number of leading
zeros grows like the inverse function of the factorial, and the
single-digit cases are by now \emph{completely} characterized:
they are exactly the $n=(i+1)!/d-1$ with $1\le d\le i$, exactly
$i$ of them per position, and the predictions $n=29,39,59,119$ as
well as $143,179,239,359,719$ all came true
(Appendix~\ref{ch:anhang-beweise}). A complete theory of the
remaining digit patterns (distribution of the number of nonzero
digits, position of the first maximal digit) is still missing.

\paragraph{The fine structure of interchange equality.} The
positivity of the interchange defect and the characterization of
the equality cases are proved (Appendix~\ref{ch:anhang-beweise});
the counting sequence $2$, $5$, $13$, $23$, $61$, $111$, $274$,
$564,\ldots$ obeys a semi-closed formula over the sharp quotients
$q'$ from the positivity proof, and its complete digit unfolding
has meanwhile been proved: per level two inequalities, a growing
excess, one descent (Appendix~\ref{ch:anhang-beweise}). The
asymptotics of the zero-count proportion have meanwhile been
settled from above: indifference to the order of growth is
asymptotically extremely rare -- the proportion falls at least
like $e^{-(1+o(1))\,n\ln n}$, because sharpness chokes off the
digits level by level. What remains open is the matching lower
bound.

\paragraph{The upper side of the depth law.} The lower side is by
now completely proved -- jump positions, asymptotics and, as the
last piece during the final revision, the exact equality: the
maximal depth in shell $m$ is \emph{exactly} the shell gap, and an
always serviceable witness is the terminal representation of the
value of the maximal polynomial (for shell 3: the house-spirit
number $(6,6)$).
All that remains open is the upper side: how fast does the
\emph{maximum} of the commutator grow? Here caution is called for,
and methodically well-founded caution at that: only the first jump
is exhaustively secured ($+2$ from the ninth world onwards). All
later thresholds are merely bounded from above ($+3$ from the
fourteenth at the latest, $+4$ from the twentieth at the latest,
$+5$ from the twenty-sixth at the latest) -- every deepening of
the witness search has so far pushed them forwards, and two
premature formula hypotheses died at exactly this front, one of
them before it was even written down.

\section{The structural questions}

\paragraph{Coupled arithmetics.} The side-fidelity theorem forbids
mixing the identities; \emph{within} one side, however, the
landscape is not classified. The product constructions via the
excess supply a whole family; the conservative-compact semiring is
demonstrably not one of them. Are there further coupled
structures? Is there a complete classification per side?

\paragraph{Addendum: the non-commutative mixed case is closed.}
This gap, too, fell during the final revision: if one demands, as
in every semiring, \emph{both} distributive laws, then right
distributivity mirrors the chain construction and replaces
commutativity entirely (proof in
Appendix~\ref{ch:anhang-beweise}). What remains is only the
curiosity of a one-sidedly distributive mixed system -- outside
any semiring axiomatics, and honestly more a footnote than a
question.

\paragraph{Addendum: three-sided side fidelity is proved.}
During the final revision the last case fell as well
(shape addition with fraction multiplication): the chain equation
forces the polynomial identity $2A(Y)=Y^{\underline{k}}$, which
demands the coefficient $\tfrac12$ -- but digits are integers.
Together with the two previously settled cases this gives: there
are exactly two worlds of arithmetic (value and shape), and the
fraction carries no third and mixes into neither. Details in
Appendix~\ref{ch:anhang-beweise}. So what remains open here is --
nothing at all; the question moves into
Chapter~\ref{ch:arithmetik} as a proved theorem.

\paragraph{Formalization.} The side-fidelity theorem is elementary
enough for a complete formalization in Lean; the falling
factorials are already available in Mathlib. That would be the
gold standard of assurance -- and a contribution to the
formal-methods community into the bargain.

\section{What the theory wants to be measured against}

Finally, the question behind all questions: when would Horimetrik
have \emph{failed}, when would it have \emph{succeeded}? Failed,
if an existing formalism is found in which horizon, duality and
side fidelity are mere renamings -- the literature search has not
found one, but searches prove nothing. Succeeded, if one of three
things happens: if the counting sequences of this book turn up
independently elsewhere; if the side-fidelity theorem or the depth
law is proved anew in a foreign language; or if someone who has
never read this book finds the same delight in the digits of $n!$
in its own shell as we did.

\begin{oberflaeche}
A theory is finished when it no longer generates questions. By
that standard Horimetrik is delightfully unfinished -- and that is
exactly how this stop along the way should end: not with a full
stop, but with a colon. What remains is the question of the
address: where on the map does all of this live?
\end{oberflaeche}
